% Options for packages loaded elsewhere
% Options for packages loaded elsewhere
\PassOptionsToPackage{unicode}{hyperref}
\PassOptionsToPackage{hyphens}{url}
\PassOptionsToPackage{dvipsnames,svgnames,x11names}{xcolor}
%
\documentclass[
  spanish,
  11pt,
  letterpaper,
]{article}
\usepackage{xcolor}
\usepackage[lmargin=2.54cm,rmargin=2.54cm,tmargin=2.54cm,bmargin=2.54cm]{geometry}
\usepackage{amsmath,amssymb}
\setcounter{secnumdepth}{5}
\usepackage{iftex}
\ifPDFTeX
  \usepackage[T1]{fontenc}
  \usepackage[utf8]{inputenc}
  \usepackage{textcomp} % provide euro and other symbols
\else % if luatex or xetex
  \usepackage{unicode-math} % this also loads fontspec
  \defaultfontfeatures{Scale=MatchLowercase}
  \defaultfontfeatures[\rmfamily]{Ligatures=TeX,Scale=1}
\fi
\usepackage{lmodern}
\ifPDFTeX\else
  % xetex/luatex font selection
  \setmainfont[]{Times New Roman}
\fi
% Use upquote if available, for straight quotes in verbatim environments
\IfFileExists{upquote.sty}{\usepackage{upquote}}{}
\IfFileExists{microtype.sty}{% use microtype if available
  \usepackage[]{microtype}
  \UseMicrotypeSet[protrusion]{basicmath} % disable protrusion for tt fonts
}{}
\makeatletter
\@ifundefined{KOMAClassName}{% if non-KOMA class
  \IfFileExists{parskip.sty}{%
    \usepackage{parskip}
  }{% else
    \setlength{\parindent}{0pt}
    \setlength{\parskip}{6pt plus 2pt minus 1pt}}
}{% if KOMA class
  \KOMAoptions{parskip=half}}
\makeatother
% Make \paragraph and \subparagraph free-standing
\makeatletter
\ifx\paragraph\undefined\else
  \let\oldparagraph\paragraph
  \renewcommand{\paragraph}{
    \@ifstar
      \xxxParagraphStar
      \xxxParagraphNoStar
  }
  \newcommand{\xxxParagraphStar}[1]{\oldparagraph*{#1}\mbox{}}
  \newcommand{\xxxParagraphNoStar}[1]{\oldparagraph{#1}\mbox{}}
\fi
\ifx\subparagraph\undefined\else
  \let\oldsubparagraph\subparagraph
  \renewcommand{\subparagraph}{
    \@ifstar
      \xxxSubParagraphStar
      \xxxSubParagraphNoStar
  }
  \newcommand{\xxxSubParagraphStar}[1]{\oldsubparagraph*{#1}\mbox{}}
  \newcommand{\xxxSubParagraphNoStar}[1]{\oldsubparagraph{#1}\mbox{}}
\fi
\makeatother


\usepackage{longtable,booktabs,array}
\usepackage{calc} % for calculating minipage widths
% Correct order of tables after \paragraph or \subparagraph
\usepackage{etoolbox}
\makeatletter
\patchcmd\longtable{\par}{\if@noskipsec\mbox{}\fi\par}{}{}
\makeatother
% Allow footnotes in longtable head/foot
\IfFileExists{footnotehyper.sty}{\usepackage{footnotehyper}}{\usepackage{footnote}}
\makesavenoteenv{longtable}
\usepackage{graphicx}
\makeatletter
\newsavebox\pandoc@box
\newcommand*\pandocbounded[1]{% scales image to fit in text height/width
  \sbox\pandoc@box{#1}%
  \Gscale@div\@tempa{\textheight}{\dimexpr\ht\pandoc@box+\dp\pandoc@box\relax}%
  \Gscale@div\@tempb{\linewidth}{\wd\pandoc@box}%
  \ifdim\@tempb\p@<\@tempa\p@\let\@tempa\@tempb\fi% select the smaller of both
  \ifdim\@tempa\p@<\p@\scalebox{\@tempa}{\usebox\pandoc@box}%
  \else\usebox{\pandoc@box}%
  \fi%
}
% Set default figure placement to htbp
\def\fps@figure{htbp}
\makeatother



\ifLuaTeX
\usepackage[bidi=basic]{babel}
\else
\usepackage[bidi=default]{babel}
\fi
\ifPDFTeX
\else
\babelfont{rm}[]{Times New Roman}
\fi
% get rid of language-specific shorthands (see #6817):
\let\LanguageShortHands\languageshorthands
\def\languageshorthands#1{}


\setlength{\emergencystretch}{3em} % prevent overfull lines

\providecommand{\tightlist}{%
  \setlength{\itemsep}{0pt}\setlength{\parskip}{0pt}}



 


\usepackage{indentfirst}
\setlength{\parindent}{1.27cm}
\makeatletter
\@ifpackageloaded{caption}{}{\usepackage{caption}}
\AtBeginDocument{%
\ifdefined\contentsname
  \renewcommand*\contentsname{Tabla de contenidos}
\else
  \newcommand\contentsname{Tabla de contenidos}
\fi
\ifdefined\listfigurename
  \renewcommand*\listfigurename{Listado de Figuras}
\else
  \newcommand\listfigurename{Listado de Figuras}
\fi
\ifdefined\listtablename
  \renewcommand*\listtablename{Listado de Tablas}
\else
  \newcommand\listtablename{Listado de Tablas}
\fi
\ifdefined\figurename
  \renewcommand*\figurename{Figura}
\else
  \newcommand\figurename{Figura}
\fi
\ifdefined\tablename
  \renewcommand*\tablename{Tabla}
\else
  \newcommand\tablename{Tabla}
\fi
}
\@ifpackageloaded{float}{}{\usepackage{float}}
\floatstyle{ruled}
\@ifundefined{c@chapter}{\newfloat{codelisting}{h}{lop}}{\newfloat{codelisting}{h}{lop}[chapter]}
\floatname{codelisting}{Listado}
\newcommand*\listoflistings{\listof{codelisting}{Listado de Listados}}
\makeatother
\makeatletter
\makeatother
\makeatletter
\@ifpackageloaded{caption}{}{\usepackage{caption}}
\@ifpackageloaded{subcaption}{}{\usepackage{subcaption}}
\makeatother
\usepackage{bookmark}
\IfFileExists{xurl.sty}{\usepackage{xurl}}{} % add URL line breaks if available
\urlstyle{same}
\hypersetup{
  pdftitle={Bitacora\_1},
  pdflang={es-CR},
  colorlinks=true,
  linkcolor={blue},
  filecolor={Maroon},
  citecolor={Blue},
  urlcolor={blue},
  pdfcreator={LaTeX via pandoc}}


\title{Bitacora\_1}
\author{}
\date{}
\begin{document}
\maketitle

\renewcommand*\contentsname{Índice}
{
\hypersetup{linkcolor=black}
\setcounter{tocdepth}{3}
\tableofcontents
}

\newpage

\section{Parte de planificación}\label{parte-de-planificaciuxf3n}

\subsection{Pregunta de
investigación}\label{pregunta-de-investigaciuxf3n}

¿Cómo influyen variables como la edad del asegurado, la antigüedad con
la aseguradora, el historial de siniestros y las características del
vehículo en la frecuencia y severidad de los siniestros, y qué factores
explican el incremento abrupto en la cantidad de reclamos cuando estos
superan el umbral de 900?

\subsubsection{Definición de la idea}\label{definiciuxf3n-de-la-idea}

La idea de esta investigación es comprobar si existe una relación entre
la edad del asegurado, tipo de vehículo, peso y largo del vehículo,
valor del vehículo, etc;con el número de siniestros y la severidad de
los mismos. Además de esto, se pueden predecir los siniestros que un
asegurado podría tener, categorizar a cada asegurado con un nivel de
riesgo, predecir el costo del siniestro, encontrar posibles casos de
fraude y ver en qué tipo de vehículo se da mayor cantidad de siniestros
y cuáles son los que tienen el mayor monto reclamado además de estimar
el monto del siniestro a partir de la zona.

\subsubsection{Conceptualización de la
idea}\label{conceptualizaciuxf3n-de-la-idea}

\subsubsection{Identificación de
tensiones}\label{identificaciuxf3n-de-tensiones}

Algunas tensiones que se pueden encontrar durante la realización del
trabajo son que, en los datos, puede haber errores o inconsistencias, ya
que no se conoce con exactitud cómo fueron recolectados. Además, al
tratarse de datos que provienen de una región específica, puede resultar
difícil generalizar los resultados a otros contextos. Por otro lado, al
abarcar únicamente el periodo de 2015 a 2018, no se consideran
situaciones importantes como la pandemia o la situación económica
actual, por lo que los resultados podrían no ser completamente
aplicables a la actualidad.

\subsubsection{Reformulación de la idea en modo
pregunta}\label{reformulaciuxf3n-de-la-idea-en-modo-pregunta}

\begin{enumerate}
\def\labelenumi{\arabic{enumi}.}
\item
  (Pregunta elegida) ¿Cómo influyen variables como la edad del
  asegurado, la antigüedad con la aseguradora, el historial de
  siniestros y las características del vehículo en la frecuencia y
  severidad de los siniestros, y qué factores explican el incremento
  abrupto en la cantidad de reclamos cuando el costo del siniestro
  supera aproximadamente las 900 unidades monetarias?
\item
  ¿Cómo varía la distribución de la frecuencia y severidad de los
  siniestros reportados entre los distintos perfiles de asegurados y
  condiciones del seguro, y qué factores explican dichas diferencias?
\item
  ¿Qué patrones se pueden encontrar en la ocurrencia de los siniestros
  reportados que permiten distinguir distintos niveles de riesgo entre
  los asegurados?
\item
  ¿Hasta qué punto las características físicas del vehículo como la
  potencia, el largo, el peso y otros influyen en mayor o menor medida
  en la frecuencia y severidad de los siniestros?
\end{enumerate}

Ahora, la opción 2 se descartó puesto que plantea un enfoque
principalmente descriptivo, centrado en analizar cómo varía la
distribución de la frecuencia y severidad de los siniestros entre
distintos perfiles. Sin embargo, no incorpora de manera explícita un
objetivo explicativo más profundo ni la identificación de fenómenos
específicos de interés dentro de los datos, como cambios abruptos en el
comportamiento de los reclamos. En cuanto a la opción 3, se descartó ya
que se enfoca principalmente en la identificación de patrones y en la
clasificación de asegurados según niveles de riesgo. No obstante, esta
formulación no aborda lo que se pensó inicialmente con el proyecto, que
es directamente la relación entre variables específicas y la frecuencia
y severidad de los siniestros. Dado esto, no se podría abordar la base
de datos de la manera deseada. La opción 4 se descartó pues está
contenida de forma implícita en la pregunta de investigación
seleccionada, además de que solamente contempla características físicas
del vehículo y no otros factores importantes como la edad del asegurado
o la zona de circulación.

\subsubsection{Argumentación de la
pregunta}\label{argumentaciuxf3n-de-la-pregunta}

\paragraph{Pregunta \#1 ¿Cómo influyen variables como la edad del
asegurado, la antigüedad con la aseguradora, el historial de siniestros
y las características del vehículo en la frecuencia y severidad de los
siniestros, y qué factores explican el incremento abrupto en la cantidad
de reclamos cuando el costo del siniestro supera aproximadamente las 900
unidades
monetarias?}\label{pregunta-1-cuxf3mo-influyen-variables-como-la-edad-del-asegurado-la-antiguxfcedad-con-la-aseguradora-el-historial-de-siniestros-y-las-caracteruxedsticas-del-vehuxedculo-en-la-frecuencia-y-severidad-de-los-siniestros-y-quuxe9-factores-explican-el-incremento-abrupto-en-la-cantidad-de-reclamos-cuando-el-costo-del-siniestro-supera-aproximadamente-las-900-unidades-monetarias}

\textbf{Contraargumentos}

\textbf{Lógica:} La pregunta combina múltiples variables explicativas y
dos dimensiones distintas del análisis (frecuencia y severidad), lo que
puede dificultar la identificación clara de relaciones específicas.
Además, el umbral de 900 unidades monetarias surge de la observación de
un gráfico específico durante el análisis exploratorio, por lo que
corresponde a una evidencia empírica inicial y podría no ser
generalizable ni representar un punto de cambio real sin un análisis
estadístico más riguroso.

\textbf{Ética:} El uso de variables como la edad o el historial de
siniestros puede implicar riesgos de segmentación injusta si no se
justifica adecuadamente su inclusión, especialmente si se utilizan para
la toma de decisiones que afecten a los asegurados.

\textbf{Emocional:} El énfasis en factores de riesgo y en el incremento
de reclamos puede percibirse como una visión reduccionista del
asegurado, sin considerar el contexto individual detrás de cada
siniestro, lo que podría generar una percepción de trato impersonal.

\textbf{Argumentos}

\textbf{Lógica:} Mediante el uso de modelos estadísticos adecuados y
modelos separados para frecuencia (por ejemplo, Poisson) y severidad
(por ejemplo, Gamma), se puede identificar de manera clara cómo factores
como la edad, la antigüedad y las características del vehículo influyen
en los siniestros. Además, el análisis del comportamiento alrededor del
umbral de 900 unidades monetarias permite explorar posibles patrones en
la distribución de los costos.

\textbf{Ética:} A través de una selección responsable de variables y una
evaluación de posibles sesgos en los datos, se puede asegurar que el
análisis no promueva discriminación injusta, sino que se base en
criterios transparentes. Asimismo, el uso adecuado de la información
permite generar conclusiones que contribuyan a mejorar la equidad en la
toma de decisiones dentro del ámbito asegurador.

\textbf{Emocional:} La interpretación de los resultados considerando el
contexto de los asegurados y comunicándolos de manera clara puede
generar mayor confianza en el análisis. Además, comprender cómo y por
qué ocurren los siniestros permite una mejor conexión con la realidad de
los usuarios, evitando percepciones de arbitrariedad o trato impersonal.

\textbf{Conclusiones} La propuesta para analizar la influencia de
variables como la edad del asegurado, el historial de siniestros y las
características del vehículo en la frecuencia y severidad de los
siniestros debe integrar un enfoque estadístico riguroso que permita
identificar patrones relevantes (Lógica), junto con una consideración
ética en el uso de variables y la posible presencia de sesgos en los
datos (Ética), y una interpretación de los resultados que conecte con la
realidad de los asegurados y evite percepciones de trato impersonal
(Emocional). Asimismo, el análisis del umbral de 900 unidades monetarias
debe entenderse como una observación empírica inicial que requiere
validación adicional. Esta combinación permite generar conclusiones más
completas, responsables y útiles para la toma de decisiones en el ámbito
asegurador.

\paragraph{Pregunta \#2 ¿Cómo varía la distribución de la frecuencia y
severidad de los siniestros reportados entre los distintos perfiles de
asegurados y condiciones del seguro, y qué factores explican dichas
diferencias?}\label{pregunta-2-cuxf3mo-varuxeda-la-distribuciuxf3n-de-la-frecuencia-y-severidad-de-los-siniestros-reportados-entre-los-distintos-perfiles-de-asegurados-y-condiciones-del-seguro-y-quuxe9-factores-explican-dichas-diferencias}

\textbf{Contraargumentos}

\textbf{Lógica:} La pregunta toma en cuenta la frecuencia y la
severidad, las cuales son conceptos distintos, así como conceptos
amplios como ``perfiles de asegurados'' y ``condiciones del seguro'', lo
que puede dificultar un análisis claro y específico.

\textbf{Ética:} El análisis de diferencias entre perfiles de asegurados
podría implicar el uso de variables sensibles (como edad u otras
características personales), lo que podría derivar en segmentaciones
injustas o en la reproducción de sesgos presentes en los datos.

\textbf{Emocional:} El enfoque en la diferenciación entre perfiles puede
percibirse como reduccionista, al no considerar las realidades
individuales detrás de los datos, lo que podría generar incomodidad o
una percepción de trato impersonal hacia los asegurados.

\textbf{Argumentos}

\textbf{Lógica:} Mediante el uso de técnicas estadísticas adecuadas,
como modelos separados para frecuencia (por ejemplo, Poisson) y
severidad (por ejemplo, Gamma), así como la correcta definición de
variables como ``perfiles de asegurados'' y ``condiciones del seguro'',
se puede obtener un análisis más preciso y bien estructurado que capture
las diferencias en los siniestros.

(Importante para bibliografía: 2.2.3.2 Poisson Distribution, 3.2.1 Gamma
Distribution
https://openacttexts.github.io/Loss-Data-Analytics/ChapFrequency-Modeling.html)

\textbf{Ética:} Utilizando una selección de variables y una revisión de
posibles sesgos en los datos, se puede garantizar que el análisis no
promueva discriminación injusta, sino que esté basada en criterios
estadísticos transparentes y responsables.

\textbf{Emocional:} La interpretación de los resultados considerando el
contexto de los asegurados y comunicándolos de forma clara y
comprensible puede generar mayor confianza y conexión, evitando
percepciones de trato impersonal o injusto.

\textbf{Conclusiones} La propuesta para analizar cómo varía la
distribución de la frecuencia y severidad de los siniestros entre
distintos perfiles de asegurados y condiciones del seguro debe basarse
en una adecuada delimitación de variables y en el uso de modelos
estadísticos apropiados que permitan capturar dichas diferencias de
manera clara (Lógica), considerando al mismo tiempo los posibles riesgos
de sesgo y discriminación asociados al uso de ciertas variables (Ética),
y procurando una comunicación de resultados que sea comprensible y
cercana a la realidad de los asegurados (Emocional). Esta integración
favorece un análisis más equilibrado y transparente, que no solo
describe diferencias, sino que también contribuye a una mejor
comprensión y aplicación de los resultados en contextos reales.

\paragraph{Pregunta 3 ¿Qué patrones se pueden encontrar en la ocurrencia
de los siniestros reportados que permiten distinguir distintos niveles
de riesgo entre los
asegurados?}\label{pregunta-3-quuxe9-patrones-se-pueden-encontrar-en-la-ocurrencia-de-los-siniestros-reportados-que-permiten-distinguir-distintos-niveles-de-riesgo-entre-los-asegurados}

\textbf{Contraargumentos}

\textbf{Lógica:} La pregunta asume que se van a encontrar patrones en
los datos, cuando aún no se ha realizado el análisis, lo que podría
sesgar la interpretación de los resultados o llevar a conclusiones
forzadas.

\textbf{Ética:} La búsqueda de ``patrones'' puede derivar en
segmentaciones de los asegurados que, aunque sean estadísticamente
válidas, podrían resultar en prácticas discriminatorias o injustas si no
se analizan con el debido cuidado.

\textbf{Emocional:} El enfoque en la identificación de niveles de riesgo
puede percibirse como una clasificación rígida de los asegurados, sin
considerar sus circunstancias individuales, lo que podría generar
incomodidad o una sensación de etiquetamiento.

\subsubsection{Argumentación a través de
datos}\label{argumentaciuxf3n-a-travuxe9s-de-datos}

\paragraph{Fuente de información}\label{fuente-de-informaciuxf3n}

Los datos fueron obtenidos de la Universidad de Valencia a través de la
página Mendeley Data en la cual se descargaron dos archivos tipo csv y
uno tipo excel, el primer archivo se titula ``Descriptive of the
variables.xlsx'' en el cual se encuentra la descripción de cada una de
las variables/columnas de los datos, el segundo archivo se titula
``Motor vehicle insurance data.csv'', este archivo posee la información
de las pólizas de seguros y el tercer archivo se titula ``sample type
claim.csv'' el cual posee información detallada de los siniestros, donde
se indica el ID, el monto total de reclamos en un año, el costo de una
parte de los reclamos y la categorización del tipo de reclamo.

\paragraph{Contexto temporal y espacial de los
datos}\label{contexto-temporal-y-espacial-de-los-datos}

Los datos contemplan la información de pólizas de seguros de autos desde
1980 hasta el 2017 según la fecha de inicio de pólizas y del 2015 al
2019 según las fechas de renovaciones; además de esto, hay información
de los siniestros para los años del 2016 al 2018. El contexto espacial
corresponde a la ciudad de Valencia, España.

\paragraph{Facilidad de obtener la
infromación}\label{facilidad-de-obtener-la-infromaciuxf3n}

Como se mencionó anteriormente, los datos se extrajeron de la página
``Mendeley Data'', para descargar los datos no se requiere solicitarlos,
cualquier persona puede descargarlos gratuitamente; sin embargo, algunas
columnas como ``length'' y ``weigth'' poseen algunos NA, además, las
columnas que poseen fechas, como ``Date\_start\_contract'' y otras,
vienen como tipo character y no como tipo date por lo que se debe hacer
la transformación de tipo character a tipo date.

\paragraph{Población de estudio}\label{poblaciuxf3n-de-estudio}

La población estudio de esta investigación corresponde a los asegurados
que posee al menos una póliza de seguro de vehículos en la ciudad de
Valencia; donde se contemplan motocicletas, furgonetas, turismos
(automóviles de 4 a 5 asientos) y vehículos agrícolas; ya sean de zonas
rurales o urbanas; con información de los siniestros que van del 2016 al
2018.

\paragraph{Muestra observada}\label{muestra-observada}

La muestra observada posee un total de 105555 observaciones, de las
cuales 53502 corresponden a identificaciones únicas en las que cada una
es un asegurado distinto; además se tienen 30 columnas en total, además
posee asegurados con y sin siniestros.

\paragraph{Unidad estadística o
individuos}\label{unidad-estaduxedstica-o-individuos}

La unidad estadística es póliza por año, donde cada fila de los datos
corresponde a una anualidad del contrato del asegurado; esto hace que un
mismo asegurado pueda aparecer más de una vez en las filas, por eso se
pasó de 105555 observaciones totales a 53502 únicas; como puede aparecer
múltiples veces, da la posibilidad de ver la evolución de la póliza a lo
largo del tiempo. Mientras que en el archivo ``sample type claim.csv'',
la unidad estadística es el tipo de siniestro para cada individuo.

\paragraph{Descripción de las variales de la
tabla}\label{descripciuxf3n-de-las-variales-de-la-tabla}

\subparagraph{Motor vehicle insurance
data}\label{motor-vehicle-insurance-data}

\begin{itemize}
\item
  \textbf{ID} (Identificador): Corresponde al número de identificación
  interna asignada a cada contrato anual de un asegurado
\item
  \textbf{Date\_start\_contract} (Fecha de inicio del contrato): Fecha
  de inicio del contrato del asegurado, tiene formato día/mes/año
\item
  \textbf{Date\_last\_renewal} (Fecha de la última renovación): Fecha de
  la última renovación del contrato, tiene formato día/mes/año
\item
  \textbf{Date\_next\_renewal} (Fecha de la próxima renovación): Fecha
  de la próxima renovación del contrato, tiene formato día/mes/año
\item
  \textbf{Distribution\_channel} (Canal de distribución): Clasifica el
  canal a través del cual se contrató la póliza. 0 si fue un Agente y 1
  para Corredores de seguros.
\item
  \textbf{Date\_birth} (Fecha de nacimiento): Fecha de nacimiento del
  asegurado declarada en la póliza, tiene formato día/mes/año.
\item
  \textbf{Date\_driving\_licence} (Fecha de la licencia de conducir):
  Fecha de expedición del permiso de conducir de la persona asegurada,
  tiene formato día/mes/año.
\item
  \textbf{Seniority} (Antigüedad): Número total de años que el asegurado
  ha estado asociado con la entidad aseguradora.
\item
  \textbf{Policies\_in\_force} (Pólizas en vigencia): Número total de
  pólizas que el asegurado posee con la entidad aseguradora en el
  periodo de referencia.
\item
  \textbf{Max\_policies} (Máximo de pólizas): Número máximo de pólizas
  que el asegurado ha tenido alguna vez en vigor con la entidad
  aseguradora.
\item
  \textbf{Max\_products} (Productos máximos): Número máximo de productos
  que el asegurado ha mantenido simultáneamente en un momento dado.
\item
  \textbf{Lapse} (Cancelación): Número de pólizas que el cliente ha
  cancelado o se le han cancelado por impago en el año de vencimiento
  actual, excluyendo aquellas que han sido reemplazadas por otra póliza.
\item
  \textbf{Date\_lapse} (Fecha de cancelación): Fecha de cancelación;
  fecha de terminación del contrato, tiene formato día/mes/año
\item
  \textbf{Payment} (Método de pago): Último método de pago de la póliza
  de referencia. El 1 representa un proceso administrativo semestral y 0
  indica un método de pago anual.
\item
  \textbf{Premium} (Prima): Monto de la prima neta asociada a la póliza
  durante el año en curso.
\item
  \textbf{Cost\_claims\_year} (Costo de siniestro del año): Costo total
  de los siniestros de la póliza de seguro durante el año en curso.
\item
  \textbf{N\_claims\_year} (Número de siniestro en el año): Número total
  de siniestros ocurridos para la póliza de seguro durante el año en
  curso.
\item
  \textbf{N\_claims\_history} (Número de siniestros histórico): Número
  total de siniestros declarados durante toda la vigencia de la póliza
  de seguro.
\item
  \textbf{R\_Claims\_history} (Ratio de siniestros histórico): Ratio del
  número de siniestros declarados para la póliza específica con respecto
  a la duración total (en años completos) de la póliza en vigor.
\item
  \textbf{Type\_risk} (Tipo de riesgo): Tipo de riesgo asociado a la
  póliza. Cada valor corresponde a un tipo de riesgo específico: 1 para
  motocicletas, 2 para furgonetas, 3 para turismos y 4 para vehículos
  agrícolas.
\item
  \textbf{Area} (Zona): Representa la zona de circulación del vehículo.
  0 para rural y 1 para urbano (más de 30.000 habitantes) en términos de
  condiciones de tráfico.
\item
  \textbf{Second\_driver} (Segundo conductor): 1 si hay varios
  conductores habituales declarados, o 0 si solo se declara un
  conductor.
\item
  \textbf{Year\_matriculation} (Año de matriculación): Año de
  matriculación del vehículo.
\item
  \textbf{Power} (Potencia): Potencia del vehículo medida en caballos de
  fuerza.
\item
  \textbf{Cylinder\_capacity} (Capacidad de cilindrada): Cilindrada del
  vehículo.
\item
  \textbf{Value\_vehicle} (Valor del vehículo): Valor de mercado del
  vehículo al 31 de diciembre del 2019.
\item
  \textbf{N\_doors} (Número de puertas): Número total de puertas del
  vehículo.
\item
  \textbf{Type\_fuel} (Tipo de combustible): Tipo específico de
  combustible utilizado para el vehículo. P para Gasolina y D para
  Diésel.
\item
  \textbf{Length} (Longitud): Longitud en metros del vehículo.
\item
  \textbf{Weight} (Peso): Peso en kilogramos del vehículo.
\end{itemize}

\subparagraph{Sample type claim}\label{sample-type-claim}

\begin{itemize}
\item
  \textbf{ID} (Identificador): Mismo identificador que en el archivo
  principal. Permite vincular cada siniestro con la póliza asociada.
\item
  \textbf{Cost\_claims\_year} (Costo de siniestros del año): Costo total
  de los siniestros de la póliza durante el año en curso.
\item
  \textbf{Cost\_claims\_by\_type} (Costo de siniestros por tipo): Costo
  asociado a un tipo específico de reclamo dentro del siniestro.
\item
  \textbf{Claims\_type} (Tipo de siniestro): Clasificación del tipo de
  reclamo como: asistencia en viaje, rotura de lunas, reclamación,
  negligencia, robo, incendio, lesiones, etc.
\end{itemize}

\subsection{Revisión bibliográfica}\label{revisiuxf3n-bibliogruxe1fica}

\subsubsection{Búsqueda de
bibliografía}\label{buxfasqueda-de-bibliografuxeda}

\subsubsection{Construcción de fichas de
literatura}\label{construcciuxf3n-de-fichas-de-literatura}

\subsection{Construcción de la UVE de
Gowin}\label{construcciuxf3n-de-la-uve-de-gowin}

\subsubsection{Conceptos básicos}\label{conceptos-buxe1sicos}

\section{Parte de escritura}\label{parte-de-escritura}

\section{Arco narrativo del proyecto}\label{arco-narrativo-del-proyecto}

\section{Diario de aprendizaje}\label{diario-de-aprendizaje}




\end{document}
